\documentclass[11pt]{article}

\usepackage[margin=1.05in]{geometry}
\usepackage{amsmath,amssymb,amsthm,mathtools}
\usepackage{enumitem}
\usepackage[hidelinks]{hyperref}
\usepackage{microtype}

\newtheorem{theorem}{Theorem}[section]
\newtheorem{lemma}[theorem]{Lemma}
\newtheorem{proposition}[theorem]{Proposition}
\newtheorem{corollary}[theorem]{Corollary}
\newtheorem{definition}[theorem]{Definition}
\newtheorem{remark}[theorem]{Remark}
\newtheorem{statement}[theorem]{Statement}
\newtheorem{conjecture}{Conjecture}
\newtheorem{question}{Question}
\renewcommand{\thequestion}{\Alph{question}}

% A theorem-like environment whose full heading is given as an argument.
\newcommand{\namedname}{}
\newtheorem*{namedinner}{\namedname}
\newenvironment{named}[1]{\renewcommand{\namedname}{#1}\begin{namedinner}}{\end{namedinner}}

\newcommand{\Z}{\mathbb Z}
\newcommand{\N}{\mathbb N}
\newcommand{\R}{\mathbb R}
\newcommand{\C}{\mathbb C}
\newcommand{\Q}{\mathbb Q}
\newcommand{\PP}{\mathbb P}
\newcommand{\GI}{\mathbb Z[i]}
\newcommand{\Norm}{\operatorname{N}}
\newcommand{\ord}{\operatorname{ord}}
\newcommand{\dist}{\operatorname{dist}}
\newcommand{\floor}[1]{\left\lfloor #1 \right\rfloor}
\newcommand{\ceil}[1]{\left\lceil #1 \right\rceil}
\newcommand{\conj}[1]{\overline{#1}}
\newcommand{\E}{\mathcal E}
\newcommand{\M}{\mathcal M}
\newcommand{\LL}{\mathcal L}
\newcommand{\Mbar}{\overline M_{0,8}}
\newcommand{\Ltw}[1]{\overline L{}^{\rm tw}_{#1}}
\newcommand{\Bal}{\operatorname{Bal}}
\newcommand{\gain}{\operatorname{gain}}
\newcommand{\mult}{\operatorname{mult}}
% Record paths in typewriter type.  The second definition (the one TeX uses)
% allows line breaks after "/" and "-"; the first is what pandoc sees.
% Arguments must not contain spaces or braces.
\newcommand{\rec}[1]{\texttt{#1}}
\makeatletter
\def\rec@scan#1{\ifx#1\relax\else#1\ifx#1/\allowbreak\fi
  \ifx#1-\allowbreak\fi\expandafter\rec@scan\fi}
\expandafter\def\csname rec\endcsname#1{\texttt{\rec@scan#1\relax}}
\makeatother

\title{Lattice Points on Critical Arcs:\\ Three Diophantine Questions}
\author{Craig Chen}
\date{}

\begin{document}
\maketitle

\begin{abstract}
Let \(N=R^2\in\N\) and let \(\E_R=\{z\in\GI:|z|^2=N\}\).  The critical case
of a conjecture of Cilleruelo and Granville asserts that an arc of length
\(C\sqrt R\) on \(|z|=R\), in any position, contains a number of points of
\(\E_R\) bounded in terms of \(C\) alone.  For every large \(R\) the bound
\((\frac23+\varepsilon)\log R/\log\log R\) is proved, no bound
\(o(\log R/\log\log R)\) is known, and ``fake circles'' show that several
classes of local inputs cannot prove one.  The only standard-conjecture
routes to the uniform bound pass through Vojta's Main Conjecture for one
fixed variety.  This note records three questions isolated in a review of
these routes, each with its status and with the known partial results and
barriers: Question~A, a statement \((L_8)\) about integer \(2\times8\) matrices
(closely related to, and implied by, a statement about rational points of
\(\Mbar\) with balanced boundary contacts; the two are not known to be
equivalent), together with a lattice form that every proof must cover;
Question~B, a statement \((L_5^\circ)\) about five lattice points on a circle,
which is an approximation-constant question at the identity of the twisted
Losev--Manin \(4\)-fold, with a split calibration related to Erd\H{o}s Problem
\#887; and Question~C, whether a non-unit pencil on \(\Mbar\) can have
positive gain.  We also list what any proof must avoid.  All statements are
taken from a review record whose files we cite by path.
\end{abstract}

\tableofcontents

\section{The problem and what is known}\label{sec:problem}

\subsection{Sources and status labels}

Every statement below is taken from the review record in the directory
\rec{docs/claude\_review/} of the repository
\rec{github.com/craigxchen/Jarnik}, branch
\rec{claude/upbeat-einstein-p0peju}.  Each statement is the
referee-corrected or killer-corrected version, and each carries a status label
and a file pointer.  File paths are relative to \rec{docs/claude\_review/}.
The labels are as follows.
\begin{itemize}[leftmargin=2em]
  \item \emph{Proved} means checked by an independent adversarial referee.
  \item \emph{Proved (killer-checked)} applies to round~20 only: the result was
    checked by an adversarial ``killer'' but not by a separate referee.
  \item \emph{Proved by the killer} marks the killer's own new results
    (round~20), which no second party has checked.
  \item \emph{Conditional} means proved from a named hypothesis.
  \item \emph{Exact} means a finite exact computation.
  \item \emph{Heuristic} means a count that no proof uses.
  \item \emph{Unrefereed} means checked by its author only.
  \item \emph{Open} means neither proved nor refuted.
\end{itemize}
Theorem-like environments carry the status and the source in brackets.  A
number written next to a file path (for example ``\rec{round20/gauss.md}
Thm~4.3'') is that file's own numbering.  It is unrelated to the numbering of
this note, which is used only through cross-references.

\subsection{The problem}

Let \(N=R^2\in\N\) and \(\E_R=\{z\in\GI:|z|^2=N\}\).  For \(C>0\), let
\(M(C)\) be the supremum, over all \(R\) and over all arcs \(I\) of \(|z|=R\)
of length \(\le C\sqrt R\) \emph{in any position}, of \(\#(I\cap\E_R)\), and
let \(M_\infty(C)\) be the corresponding \(\limsup\) as \(R\to\infty\)
(\rec{overview/overview.md} \S1.1).

\begin{conjecture}[Cilleruelo--Granville, critical case; \cite{CillerueloGranville2007,CillerueloGranville2009}]\label{conj:CG}
\(M(C)<\infty\) for every \(C>0\).
\end{conjecture}

The scale \(\sqrt R\) is critical because the sagitta \(L^2/8R\) of an arc of
length \(L\) becomes of order \(1\) there.  Arcs that touch an axis point are
trivial, so the whole difficulty is uniformity in the position of the arc
(\rec{overview/overview.md} \S1.1).

\subsection{The \texorpdfstring{\((\frac23+\varepsilon)\log R/\log\log R\)}{(2/3+eps) log R/loglog R} theorem}

For every \(R\) the literature bound is \(\ll_C\log R\)
(Cilleruelo--C\'ordoba~\cite{CillerueloCordoba1992}; Bourgain--Rudnick).

\begin{theorem}[proved, two referees; \rec{two\_thirds/proof.md} Thm 7.5, Cor 8.1]\label{thm:twothirds}
Let \(M\ge2\) distinct lattice points lie on an arc of length \(\le C\sqrt R\).
Then
\[
  3M\log M\le\log R^2+(14+4\log^+C)\,M .
\]
Hence for every \(C,\varepsilon>0\) there is an explicit \(R_0(C,\varepsilon)\)
such that \(M\le(\frac23+\varepsilon)\log R/\log\log R\) for all \(R\ge R_0\).
\end{theorem}

\paragraph{Idea of the proof.}
The proof compares an archimedean upper bound with an arithmetic lower bound
for the product of the reduced chords.  After dividing by
\(\gcd(z_1,\dots,z_M)\) one obtains the \emph{normalised} cluster: \(N\) is
odd, all \(p\mid N\) split, and
\(z_j=\varepsilon_j\prod_p\pi_p^{a_{jp}}\conj\pi_p^{\,e_p-a_{jp}}\).  Put
\(c_{ij}=(z_i-z_j)/\gcd(z_i,z_j)\) (the reduced chords) and
\(S_{p\tau}=\#\{i:a_{ip}<\tau\}\) for \(1\le\tau\le e_p\).  The link is an exact
cut identity (\rec{overview/overview.md} \S2.2, Step~2):
\[
  \prod_{i<j}|z_i-z_j|^8
  =N^{M(M-2)}\prod_{p\mid N}p^{\sum_\tau(2S_{p\tau}-M)^2}\prod_{i<j}|c_{ij}|^8 .
\]
The arithmetic side comes from pigeonhole collisions at primes
\(q\lesssim M\), together with a chain lemma on exponent levels.  The factor
\(\log M\) is the mass \(\sum_{q\le M}\log q/q\).

No bound \(o(\log R/\log\log R)\) is known for every \(R\), either
unconditionally or from any standard conjecture.  Explicit examples give
\(M(C)\ge10\) for \(C\ge8.2394\) (one explicit \(10\)-point cluster) and
\(M_\infty(C)\ge8\) for \(C>5.89771\) (a family)
(\rec{growth/families.md}, \rec{growth/referee\_families.md}).

\subsection{The fake-circle barrier}\label{sec:fakes}

Why is \(\log R/\log\log R\) a barrier for local inputs
(\rec{overview/overview.md} \S\S5.1, 5.4)?  Divide a cluster by its Gaussian
gcd.  Then \(N=\prod_jp_j^{e_j}\) with \(p_j\equiv1\pmod 4\), and
\(z_x=\varepsilon_x\prod_j\pi_j^{a_{xj}}\conj\pi_j^{\,e_j-a_{xj}}\).  The arc
condition becomes a linear system in the angles \(\phi_j=\arg\pi_j\): for all
\(x\), \(\langle2a_x-e,\phi\rangle\in\theta_0+[0,\Delta]+\frac\pi2\Z\), with
\(\Delta\asymp R^{-1/2}\).  A \emph{fake circle} (angle model) keeps the
genuine primes and the genuine exponent profile but uses an arbitrary
\emph{real} \(\phi\).  That \(\phi\) must solve the system and satisfy every
monomial gap
\(\dist(\langle v,\phi\rangle,\frac\pi2\Z)\ge\arcsin e^{-w(v)/2}\), where
\(w(v)=\sum|v_j|\log p_j\); these gaps hold on actual circles.  Since
\(W=\log N=cM\log M\) corresponds to \(M\approx(2/c)\log R/\log\log R\), an
input class that admits fakes with \(W=O(M\log M)\) cannot prove
\(M=o(\log R/\log\log R)\).  The record builds such fakes for several classes.
\begin{itemize}[leftmargin=2em]
  \item Skew-Hadamard (Paley) profiles give fakes at
    \(W\le(8+o(1))M\log M+8M\log(1/C)\).  These satisfy every pair, cut,
    monomial-gap and rational-character inequality (\rec{round8/lw.md}
    Thm~6.1; \emph{proved}).
  \item Residue designs extend this to point-level residues at moduli
    \(\le M^A\) with \(A\) fixed in advance (\rec{round5/construct.md},
    \rec{round5/lower.md}).  Referee scope limits: all primes of \(N\) must
    exceed the modulus bound \(X\), and the twisted extension at \(X=M^A\),
    \(A>1\), is open (\rec{round5/README.md}; \rec{overview/overview.md}
    \S5.2 item~2).
  \item Grafts cap robust averaged analytic inputs
    (\rec{round10/analytic.md} Thm~B; \rec{round20/composite.md}).
\end{itemize}
Actual clusters satisfy \(W\ge3M\log M-O(M)\).  Not everything is capped.
P3-all\(^*\), the sharp-constant large sieves and the low-height window
\(P_+(N')\le M^{2-\delta}\) remain uncapped; here \(N'\) is the normalised
norm and \(P_+(N')\) its largest prime factor (\rec{round21/README.md}, status
\emph{open}), and the residue caps are not uniform in \(A\); see
\rec{overview/overview.md} \S5.5.  The record's synthesis is a judge's
assessment, \emph{unrefereed} (\rec{round17/README.md} \S7.1;
\rec{overview/overview.md} \S7.6).  It says a new mechanism must be global,
multi-place and tied to the actual arithmetic of \(N\).  The only statements
of that type in the record are Vojta-type, and their sharpest arithmetic forms
are Questions~\ref{q:A} and~\ref{q:B} below.

\subsection{What is conditional: Vojta}

The only standard-conjecture routes give the uniform bound directly.  They use
Vojta's Main Conjecture for \emph{one fixed} variety:
\begin{itemize}[leftmargin=2em]
  \item \(D=0\) on a resolution of \(\mathrm{Bl}_YX_7\), where \(X_7\) is the
    equal-norm variety and \(Y\) its diagonal line
    (\rec{growth/conditional.md} Thm~5.1);
  \item \(D=0\) on \(\Mbar\) (Section~\ref{sec:A});
  \item \(D=E\), \(S=\{\infty\}\) on a \(4\)-fold (Section~\ref{sec:B}).
\end{itemize}

\section{Question A: \texorpdfstring{\((L_8)\)}{(L8)}, a statement about integer \texorpdfstring{\(2\times8\)}{2x8} matrices}\label{sec:A}

\subsection{The statement}

The data are as follows (\rec{growth/ptolemy.md} \S7;
\rec{overview/overview.md} \S7.1).
\begin{itemize}[leftmargin=2em]
  \item A matrix \(P=(P_0|\cdots|P_7)\in\Z^{2\times8}\) with primitive
    columns, \(P_0=(1,0)^T\) and \(P_j=(X_j,Y_j)^T\).
  \item Positive integers \(n_T\) for \(\emptyset\ne T\subseteq[7]\), such
    that for all \(0\le i<j\le7\)
    \[
      \det(P_i,P_j)=t_{ij}\prod_{T\ni i,j}n_T,\qquad
      t_{ij}\in\Z\setminus\{0\}\quad(t_{0j}=Y_j).
    \]
  \item The \emph{chain condition}: at each prime \(p\), the sets \(T\) with
    \(p\mid n_T\) form two chains under inclusion, and the unions of the two
    chains are disjoint.
\end{itemize}
The \emph{visible} blocks are those with \(2\le|T|\le6\).  There are \(119\)
of them, one for each boundary divisor \(D_T\) of \(\Mbar\).  Singletons occur
in no minor.  The top block \([7]\) divides every minor with \(i,j\ge1\), but
it can be inflated by a fresh prime (\(z\mapsto nz\)) without moving \(\Phi\),
so imposing balance only on the \(119\) visible blocks is equivalent
(\rec{round18/README.md} \S1, ``Balance'').  Column~\(0\) is the
\emph{anchor}: it is normalised to \(P_0=(1,0)^T\) and sent to \(\infty\) in
Kapranov's chart \(\PP^5=\PP(\Q^7/\Q\mathbf 1)\).  The chain condition does
not depend on the anchor, but \(Z\) depends on the labelling, so a statement
with an arbitrary anchor must symmetrise \(Z\) (\rec{round18/README.md} \S1,
``Anchor'').  Put
\[
  H_\Sigma=\sum_j\log|Y_j|+\sum_{1\le i<j\le7}\log|t_{ij}| .
\]
The datum is \emph{\(\eta\)-balanced at scale \(w\)} if \(|\log n_T-w|\le\eta w\)
for all \(119\) visible \(T\).  Finally, \(\Phi\in M_{0,8}(\Q)\) is the point
given by the eight columns as points of \(\PP^1(\Q)\).

\begin{statement}[\((L_8)\); open]\label{st:L8}
There are \(c>0\), \(\eta_0>0\) and a proper Zariski-closed \(Z\subset\Mbar\)
such that every such datum with \(\Phi\notin Z\) and \(\eta\le\eta_0\)
satisfies \(H_\Sigma\ge c\,w-O(1)\).
\end{statement}

\subsection{Cleanest form and the sandwich}

The following reductions are \emph{proved} (\rec{round18/theory.md} \S1.3,
Prop~2.1; \rec{round18/census.md} Prop~1.3).  Every \(\Phi\in M_{0,8}(\Q)\) has
a primitive representative \(x\in\Z^7\), taken modulo translation, with
point~\(0\) sent to \(\infty\) and
\[
  |x_i-x_j|=\prod_{T\ \text{visible},\ T\ni i,j}d_T(\Phi).
\]
Here \(d_T=\prod_pp^{\ell_T(p)}\) is the \emph{contact integer}, the finite part
of the height \(h_{D_T}(\Phi)\), and \(\ell_T(p)\) is the number of levels \(t\)
at which \(T\) is a residue class of \(x\bmod p^t\) (\rec{round18/theory.md}
Lemma~2.2).  Let \(\Bal(\eta)\) be the set of \(\Phi\) with
\(|\log d_T-w|\le\eta w\) for all \(119\) \(T\).  The strongest natural form,
\((L_8^c)^{\rm all}\), which implies \((L_8)\), says that \emph{a rational point
of \(M_{0,8}\) off a fixed \(Z\) cannot have all \(119\) boundary contacts of
the same size \(e^{w}\), \(w\to\infty\)}; \((L_8)\) itself gives this only at
caterpillar points.

\begin{proposition}[proved; round 18, \rec{round18/theory.md} \S1.3, \rec{round18/census.md} Prop 1.3]\label{prop:sandwich}
\[
  (L_8^c)^{\rm all}\Rightarrow(L_8)\Rightarrow(L_8^f)\Rightarrow(L_8^c)^{\rm cat}\Rightarrow(L_8^{pc}).
\]
No arrow is known to reverse.
\end{proposition}

The terms are as follows.
\begin{itemize}[leftmargin=2em]
  \item \((L_8^c)^{\rm all}(\eta_1)\): \(\Bal(\eta_1)\cap\{w\ge w_1\}\subset Z\).
    It implies \((L_8)\) with \(\eta_0=\eta_1/2\), \(c=\eta_1/4\).  Conversely
    \((L_8)\) implies \((L_8^c)(\eta_0,Z)\) only for \(\Phi\) admissible from
    anchor~\(0\) (any anchor, with symmetrised \(Z\)); this is not an
    equivalence (\rec{round18/README.md} \S2.1 item~1;
    \rec{overview/overview.md} \S7.1 item~2).
  \item ``cat'' restricts this to points that are \emph{caterpillar} at every
    prime.  A compatible split system is a \emph{caterpillar} if one side of
    each split can be chosen so that the chosen sides form a chain;
    equivalently, all internal edges of its tree lie on one path.  The chain
    condition is exactly the caterpillar condition on the block system at each
    prime (\rec{round18/theory.md} Lemma~1.1; exact on all \(39208\) trees with
    \(8\) leaves).  Caveat (referee item~H): the block system equals the
    \(p\)-adic tree of \(\Phi\) only at primes dividing no residue.
  \item ``pc'' restricts it to pairwise coprime \(d_T\).
  \item \((L_8^f)\) replaces \(cw\) by any \(f(w)\to\infty\) with \(f\le cw\).
    Every proof of \((L_8^f)\) must prove \((L_8^c)^{\rm cat}\), the
    \(H_\Sigma=0\) content of \((L_8)\), so a slowly growing \(f\) is no
    shortcut; but \((L_8^f)\) is not known to be equivalent to \((L_8)\)
    (\rec{round18/referee\_theory.md} \S2.3, item~A).
\end{itemize}

\subsection{Lattice form}

The following is \emph{open}; it is a consequence of \((L_8)\), not known to be
equivalent to it (\rec{round18/theory.md} Lemma~2.2, \S5(2), as corrected by
\rec{round18/referee\_theory.md} item~G).
\begin{itemize}[leftmargin=2em]
  \item In Kapranov's model \(\Mbar\to\PP^5\): \(H\cdot\beta=31\),
    \(E_T\cdot\beta=1\), \(K=-6H+\sum_{3\le|T|\le6}(|T|-2)E_T\) and
    \(K\cdot\beta=10\).  Here \(\beta\) is the balanced class, defined by
    \(D_T\cdot\beta=1\) for all \(T\).
  \item The congruence lattice \(\Lambda(x)\) consists of the \(y\in\Z^7\) with
    \(y_i\equiv y_j\bmod p^t\) whenever \(i,j\) lie in one residue class of
    \(x\bmod p^t\) of size \(\ge3\).  It contains \(x\) and has index
    \([\Z^7:\Lambda(x)]=\prod_{3\le|T|\le6}d_T^{|T|-1}\).
  \item A balanced point has \(|x|\approx e^{31w}\).  The Minkowski scale of
    its own lattice is \(e^{49w}\), and the trivial (Liouville) bound is
    \(e^{30w}\).
\end{itemize}

\begin{statement}[\((L_8^{pc})\), lattice form; open]\label{st:L8pc}
There are \(\eta_1>0\), \(w_1\) and a proper closed \(Z\) with the following
property.  Let \(e_T\in[e^{(1-\eta_1)w},e^{(1+\eta_1)w}]\), for
\(T\subseteq[7]\) with \(3\le|T|\le6\), be pairwise coprime integers, and let
\(w\ge w_1\).  Then the lattice
\[
  \Lambda_e=\{y\in\Z^7:e_T\mid y_i-y_j\ \text{for } i,j\in T\}
\]
has no primitive vector \(y\notin\Z\mathbf 1\) such that all of the following
hold:
\begin{enumerate}[label=(\roman*)]
  \item \([y]\in M_{0,8}\setminus Z\);
  \item \(|y|\le e^{31w(1+\eta_1)}\);
  \item the pair cofactors \((y_i-y_j)/\prod_{T\supseteq\{i,j\}}e_T\) are
    pairwise coprime of size \(e^{w(1\pm\eta_1)}\).
\end{enumerate}
\end{statement}

In words, structured lattices have no primitive vector \(e^{18w}\) below their
Minkowski scale.  Nothing between the trivial bound and Minkowski is known for
these lattices.  The heuristic count gives \(98-108=-10=-K\cdot\beta\) per unit
of \(w\): \(e^{98w}\) profiles against \(e^{-108w}\) for shortness
(\emph{heuristic}).  \((L_8^{pc})\) alone does \emph{not} give the uniform
theorem; that needs the chain cases, that is, \((L_8)\) itself.

\subsection{What is known around \texorpdfstring{\((L_8)\)}{(L8)}}

\begin{theorem}[proved, refereed; \rec{growth/ptolemy.md} \S7 and Steps 0--5 of \S6.2, \rec{growth/referee\_ptolemy.md}]\label{thm:L8uniform}
\((L_8)\) implies the uniform bound: \(M(C)<\infty\) for every \(C>0\).
\end{theorem}

The proof has four steps.
\begin{enumerate}[leftmargin=2em]
  \item Subdivide the arc, keep one Gaussian unit class, and divide by its gcd.
  \item A random balanced selection (item~532 of the notes on branch
    \rec{agent/uniform-gcd-formalization}) shows that at least half of all
    ordered \(8\)-tuples have balance \(\eta\le2\sqrt2\cdot255/\sqrt{M-7}\) and
    \(H_\Sigma\le3584\,w/(M-1)\), with \(w=W/128\).
  \item At most \(c(Z)M^7\) of the \(8\)-tuples have \(\Phi\in Z\) (\rec{growth/ptolemy.md} Lemma~6.1).
  \item \((L_8)\) then bounds \(M\) independently of \(R\).
\end{enumerate}

\begin{theorem}[conditional on Vojta; Thm D, \rec{growth/ptolemy.md} \S6]\label{thm:VojtaL8}
Assume Vojta's Main Conjecture with \(D=0\) on \(\Mbar\) over \(\Q\), with
\(B=\sum_TD_T\) the boundary and \(A=K+B\), which is ample (Keel--McKernan).
Then \((L_8)\) holds with \(Z=Z_{1/43}\) and \(\eta_0=1/381\).
\end{theorem}

The key estimate is \(h_K(\Phi)\ge10w-282\eta w-60H_\Sigma-O(1)\).  The sign is
\(K\cdot\beta=2^{k-3}(k-8)+k+2\), which is \(-7\) at \(k=7\) and \(+10\) at
\(k=8\).  So \(k=8\) is the first case in which the \(D=0\) inequality says
anything at balanced points.

\begin{theorem}[proved; \rec{growth/referee\_ptolemy.md} \S3.3, \rec{threshold/L8.md} Prop C and Thm B]\label{thm:L56false}
\((L_5)\) and \((L_6)\) are false.
\end{theorem}

\((L_7)\) is open and expected false, and \((L_7)\Rightarrow(L_8)\)
(\rec{round18/census.md} Lemma~7.7).

\begin{remark}[Correction B; \rec{round18/referee\_theory.md} \S2.7]
A refutation of \((L_8)\) refutes Vojta at balanced points, but not
conversely: a refuter must produce caterpillar points.
\end{remark}

\subsection{Round-18 facts}

The following are from \rec{round18/README.md} \S\S0--3; all are
\emph{proved} and refereed unless labelled otherwise.
\begin{enumerate}[leftmargin=2em]
  \item \emph{No feasible computation} (\rec{round18/census.md} Thm~E).  Each
    prime divides at most \(5\) visible blocks, so a balanced datum needs at
    least \(24\) primes.  At \(\eta=0.1\) the class is empty for
    \(e^w<59.2\), and all \(21\) minors \(\det(P_i,P_j)\) with \(i,j\ge1\) are
    \(\ge10^{49.4}\).  Block-sampling searches return \(H_\Sigma\asymp w\) even
    at \(k=5,6\), where bounded-\(H\) points are known to be abundant
    (\emph{exact}).  So a failed search says nothing.
  \item \emph{No family mechanism.}
    \begin{itemize}
      \item Density, box and sieve arguments produce only
        \(\kappa\le-\sum\log e_T\), where \(\kappa\) is the finite part of
        \(h_K\) (\rec{round18/theory.md} Prop~2.3).
      \item Bounded-degree covering families of genus-\(\le1\) curves whose
        members \emph{each} carry infinitely many balanced small-\(H\) points
        are excluded (\rec{round18/census.md} Prop~7.5).  Correction~M1: a
        dense supply with finitely many points on each member, for example on
        the lines of Kapranov's \(\PP^5\), is \emph{not} excluded.
      \item Class-\(\beta_7\) lifts of the \((L_6)\) families are excluded
        (\rec{round18/lift.md} Thm~4.2).  Classes \(c\beta_7\) with \(c\ge2\),
        and point-level lifts, are open.
      \item Restricted theorems are known only on special loci with
        locus-dependent constants, or for a fixed prime set (Ridout:
        \(H_\Sigma\ge8(1-\eta)(1-\varepsilon)w-C(S_0,\eta,\varepsilon)\) when all
        blocks are \(S_0\)-units; \rec{round18/theory.md} Thm~2.5).  The
        constant \(C\) depends on \(S_0\), so this is useless here.
    \end{itemize}
  \item \emph{Chains carry no information on pairwise-coprime profiles}
    (\rec{round18/theory.md} Prop~2.1).  On \(\Bal^{pc}\) each prime divides
    one block, and the \(p\)-adic tree is a single edge.  So the chain
    condition, the anchoring and the residues are vacuous, and \emph{every
    proof of \((L_8)\) must prove \((L_8^{pc})\)}.  The gloss ``nothing but
    \(K\cdot\beta=10\) survives'' is not a theorem (item~G).
\end{enumerate}

\begin{question}\label{q:A}
Prove \((L_8)\), or at least its lattice form \((L_8^{pc})\), which every proof
must cover.  A proof of \((L_8)\) would be an unconditional case of \(D=0\)
Vojta on the rational \(5\)-fold \(\Mbar\), for a structured but unbounded
family of points.  Alternatively, find, for every \(\eta>0\) and every \(w_1\),
a Zariski-dense set of \(\eta\)-balanced points with \(w\ge w_1\) that are
caterpillar up to \(o(w)\) of residue.  That would refute both \((L_8)\) and
Vojta (Correction~B; \rec{round18/referee\_theory.md} \S2.7,
\rec{round18/theory.md} \S3.1).
\end{question}

\section{Question B: \texorpdfstring{\((L_5^\circ)\)}{(L5 circ)}, five lattice points on a circle}\label{sec:B}

Every theorem in this section is \emph{proved (killer-checked)} or
\emph{conditional}, except where labelled otherwise; the statements \(A_5\),
\((L_5^\circ)\) and its arc form are \emph{open}.  Sources:
\rec{round20/gauss.md} \S\S1--4, \rec{round20/kill\_gauss.md} \S\S1.2, 6--7,
and \rec{overview/overview.md} \S\S4.3, 7.1A.

\subsection{Setting}

\begin{itemize}[leftmargin=2em]
  \item Let \(\zeta_1,\dots,\zeta_k\) be a \emph{primitive} tuple of distinct
    lattice points of \(|\zeta|=R''\), meaning \(\gcd_{\GI}=1\).
  \item Its similarity class \([\zeta]=(\zeta_j/\zeta_1)_j\) is a point of the
    norm-one torus
    \(T_k=(R^1_{\Q(i)/\Q}\mathbb G_m)^k/\text{diag}\), which has dimension
    \(k-1\).  \(T_k\) is the open torus of the twisted Losev--Manin space
    \(\Ltw k\) (the permutohedral fan in \(\Z^k/\Z\mathbf 1\), with Galois
    acting by \(-1\) on cocharacters).  \([\zeta]\) is invariant under
    \(\zeta\mapsto\gamma\zeta\), \(\gamma\in\Q(i)^*\), which is why the
    exceptional set \(Z\) must live in \(T_k\).  For \(A\subsetneq[k]\)
    nonempty, \(E_A\) is the toric divisor of the ray \(\mathbf 1_A\).
  \item \emph{Unit class.} Fix the Gaussian primes \(\pi_p\) once.  Each point
    factors as \(\zeta=\varepsilon\prod_p\pi_p^{a(p)}\conj\pi_p^{\,e_p-a(p)}\)
    with a unit \(\varepsilon\in\{\pm1,\pm i\}\); points lie in one \emph{unit
    class} when their units agree.  Within one class,
    \(\zeta_a/\zeta_b=P/\conj P\) with \(P\in\GI\) conjugate-primitive
    (\rec{growth/ptolemy.md} \S1, following item~532 on branch
    \rec{agent/uniform-gcd-formalization}; \rec{docs/claude\_review/} has no
    separate formal definition).\footnote{The form \(P/\conj P\) is a
      consequence of the definition, not an equivalent definition:
      \(i=(1+i)/(1-i)\) and \(-1=i/(-i)\) are themselves of the form
      \(P/\conj P\).}
  \item Each threshold layer \((p,t)\) of the exponent profile gives a cut
    \(\{A,A^c\}\) of \([k]\).  Let \(n_A\) be the product of the \(p\) over the
    layers with that cut.  Then \(h_{E_A+E_{A^c}}([\zeta])=\log n_A\) exactly,
    and \(\sum_A\log n_A=\log R''^2\) (\rec{round20/gauss.md} Lemma~1.1).
  \item A short arc means archimedean proximity of \([\zeta]\) to the identity
    \(e\).
  \item Put \(d_{ab}=\sum_{A\text{ separates }a,b}\log n_A\), define \(t_{ab}\)
    by \(|\zeta_a-\zeta_b|=2|t_{ab}|R''e^{-d_{ab}/2}\), and put
    \(H_\Sigma=\sum_{a<b}\log|t_{ab}|\).  Here \(t_{ab}\) is a nonzero integer
    when \(\zeta_a,\zeta_b\) lie in one Gaussian unit class.  Across unit
    classes it need not be: \(|t_{ab}|\in\sqrt{\text{odd}/2}\) occurs (killer
    correction, minor~1).  The record allows either repair: restrict to one
    unit class (as below), or take \(t_{ab}:=|c_{ab}|/2\) real, with
    \(c_{ab}=(\zeta_a-\zeta_b)/\gcd(\zeta_a,\zeta_b)\), and read the conclusion
    as \(H_\Sigma\ge cw-O(1)\) (\rec{round20/kill\_gauss.md} \S7 item~1;
    \rec{round20/README.md} \S1.1 item~1).
\end{itemize}

\subsection{The statements}

\begin{statement}[\(A_k(\delta)\); open for \(k=5\)]\label{st:Ak}
There are a proper Zariski-closed \(Z\subset T_k\) and an \(R_0\) such that
every primitive \(k\)-tuple of distinct lattice points of a circle of radius
\(R''\ge R_0\), with all mutual distances \(\le R''^{1/2-\delta}\), has
\([\zeta]\in Z\).
\end{statement}

\begin{statement}[\((L_k^\circ)\); open for \(k=5\)]\label{st:Lkcirc}
There are \(c,\eta_0>0\), \(w_0\) and a proper Zariski-closed \(Z\subset T_k\)
with the following property.  Every primitive \(k\)-tuple \emph{within one
unit class} with \([\zeta]\notin Z\), \(w\ge w_0\) and
\(|\log n_A-w|\le\eta_0w\) for all \(2^{k-1}-1\) cuts has
\(H_\Sigma\ge cw\).
\end{statement}

At \(k=5\) there is also a sufficient \emph{arc form} (\emph{open}; it implies
\((L_5^\circ)\) by the argument of \rec{round20/gauss.md} Prop~3.1, and the
converse is not proved; \rec{overview/overview.md} \S7.1A).

\begin{statement}[arc form of \((L_5^\circ)\); open]\label{st:arc}
Take five primitive lattice points of a circle of radius \(R\), in one unit
class, whose Gaussian profile is \(\eta_0\)-balanced over all \(15\) cuts and
whose class lies outside \(Z\).  Then, once \(R\) is large, they do not all lie
within \(R^{7/15+c}\) of each other.
\end{statement}

The exponent \(7/15\) arises because each pair is separated by \(8\) of the
\(15\) cuts.  So \(|t_{ab}|\ge1\) (valid within one unit class) gives diameter
\(\ge2R''^{7/15}\) at exact balance, and this is exactly the Liouville bound.
For a balanced \(k\)-tuple taken from a large cluster on a critical
\(\sqrt R\)-arc, dividing by the tuple gcd moves it to the ``subcritical''
exponent \(\frac12-\frac1{2^k-2}\) in the reduced radius
(\rec{round20/gauss.md} Lemma~1.2).

\subsection{Theorems}

\begin{theorem}[proved, killer-checked; \rec{round20/gauss.md} Thm 3.4\('\), Prop 3.1]\label{thm:reduction}
Let \(k\ge3\) and \(0<\delta<1/(2^k-2)\).  Then
\[
  A_k(\delta)\Rightarrow(L_k^\circ)\Rightarrow\text{the uniform bound}.
\]
For \(\delta\le2^{-k-1}\), only a convexity averaging of the tuple gcd is
needed.  The converse \((L_5^\circ)\Rightarrow A_5\) is \emph{not} proved.
\end{theorem}

\begin{theorem}[proved, killer-checked; \rec{round20/gauss.md} Thm 4.1]\label{thm:k4false}
The \(4\)-point version is false.
\begin{itemize}
  \item Take Gaussian linear blocks \(H_T(s)=q_T(s-r_T)\), with roots \(r_T\) on
    an explicit rational \(2\)-parameter family.  These give balanced
    \(4\)-tuples with bounded residues at exponent
    \(3/7=\frac12-\frac1{14}\), dense in \(T_4\).
  \item Hence \(A_4(\delta)\) is false for every \(\delta<1/14\), and
    \((L_4^\circ)\) is false for every \(c,\eta_0,Z\).
\end{itemize}
\end{theorem}

The exact Jacobian rank was reproduced at \(6\) rational points.

\begin{remark}[proved by the killer, not independently checked; \rec{round20/kill\_gauss.md} \S3, \rec{round20/kill\_gauss\_checks/k3\_projection.py}]\label{rem:k3false}
Likewise \(A_3(\delta<1/6)\) and \((L_3^\circ)\) are false (the killer's own
projection corollary).
\end{remark}

So \(k=5\) is the first case where the reduction can work.

\begin{theorem}[conditional on Vojta; \rec{round20/gauss.md} Thm 4.2, re-derived by the killer]\label{thm:VE}
Let \(Y_k=\mathrm{Bl}_e\Ltw k\), so \(K_{Y_k}=-\pi^*\sum_\rho E_\rho+(k-2)E\).
Let \(V_E(k,\varepsilon)\) be Vojta's Main Conjecture for
\((Y_k,D=E,S=\{\infty\})\): off a proper closed \(Z_\varepsilon\),
\[
  \lambda_{E,\infty}(y)+h_{K_{Y_k}}(y)\le\varepsilon h_A(y)+C_\varepsilon .
\]
Here \(A=\pi^*L-\varepsilon'E\) is ample with \(L=\sum_\rho a_\rho E_\rho\),
\(a_\rho>0\), and \(c_A=\max_\rho a_\rho\), so that \(h_A\le c_AW'+O(1)\) with
\(W'=\log R''^2\) (\rec{round20/gauss.md} Thm~4.2 proof).
\begin{itemize}
  \item If \(k\ge5\), \(\kappa=(k-1)(1+2\delta)/4-1>0\) and
    \(\varepsilon<\kappa/c_A\), then \(V_E(k,\varepsilon)\Rightarrow A_k(\delta)\).
  \item At \(k=5\), \(\kappa=2\delta\), so \(V_E(5)\) with one
    \(\varepsilon<(1/32)/c_A\) gives the uniform bound (via \(A_5(1/64)\)),
    and \(V_E(5)\) predicts \((L_5^\circ)\).
\end{itemize}
The instance is non-vacuous (\rec{round20/gauss.md} Thm~4.2, Remark~(3)).  The
\(D=0\) form would need \(k\ge6\).
\end{theorem}

\begin{theorem}[proved, killer-checked; \rec{round20/gauss.md} Thm 4.3(ii)]\label{thm:covering}
Let \(k\ge5\) and \(\delta>0\).  Put \(B=\pi^*\sum_\rho E_\rho\), so
\(h_B=W'\).  Call the class of a primitive tuple with diameter
\(\le R''^{1/2-\delta}\) an \emph{A-point}.  For each \(d\), the curves \(C\)
with \(B\cdot C\le d\) that carry infinitely many A-points lie in a proper
closed set.
\end{theorem}

Part~(i) of the same result (\rec{round20/gauss.md} Thm~4.3(i)) was narrowed by the killer
(\rec{round20/kill\_gauss.md} \S1.2 item~4).  The density mechanism is excluded
at \(k=5\) only \emph{heuristically} (count \(e^{-w}\) per scale).

\begin{lemma}[K1, proved by the killer, not independently checked; \rec{round20/kill\_gauss.md} \S6.1]\label{lem:K1}
Every recorded method is exactly critical.  Let \(F\) be bihomogeneous of
bidegree \((a,b)\).
\begin{itemize}
  \item Then \(F(z)=g^a\conj g^{\,b}F(\zeta)\), and the forced divisor scales
    in the same way.
  \item The angular width is unchanged, and so are the pair data and the
    residues.
  \item Consequently, every forced-divisibility, Liouville, pair, cut,
    character or single-form inequality is the same before and after division
    by the gcd.  The move from \(\frac12\) to \(\frac7{15}\) is a change of
    units.
  \item Auxiliary forms are capped at ratio exactly \(1\)
    (\rec{threshold/construct.md} Cor~C3, with Thm~B), and single-form
    Lang--Waldschmidt is exactly tight.
  \item Only Vojta's \(D=E\) term and the heuristic count see the reduction.
\end{itemize}
\end{lemma}

\subsection{Approximation reading}

The following is from \rec{round20/kill\_gauss.md} \S6.2.  Measure height by
\(h_B=\log R''^2\).  An A-point satisfies
\(\dist(y,e)\le H_B(y)^{-(1/4+\delta/2)}\), approximating the \emph{rational}
point \(e\) of the \(4\)-fold \(T_5\).
\begin{itemize}[leftmargin=2em]
  \item \(\delta=0\) is counting-critical (\emph{heuristic}).
  \item Liouville excludes balanced points for \(\delta>1/30\), and arbitrary
    points for \(\delta>1/10\) (Cilleruelo--C\'ordoba).
  \item What is needed is that, for some \(\delta<1/30\), \emph{the Liouville
    exponent at \(e\) is not attained by a Zariski-dense set of balanced
    points}.
  \item At \(k=4\) this fails (Theorem~\ref{thm:k4false}).  On \(\PP^1\),
    Liouville is attained at every rational point.
  \item Roth- and Schmidt-type theorems apply but are weaker.
    McKinnon--Roth~\cite{McKinnonRoth2015} Thm~5.1 gives, off a proper closed
    set, \(\alpha^{\rm ess}_x\ge\beta_x(L)\), where
    \(\beta_x(L)=\int_0^{\tau}\mathrm{Vol}(\pi^*L-\gamma E)/\mathrm{Vol}(L)\,d\gamma<\tau_x(L)\)
    (ibid.\ p.~25).  Their Thm~6.3 gives \(\alpha_x\ge\varepsilon_x/2\), strict
    at rational \(x\).  At \(e\) (on \(X'\)), \(\beta_e(-K)\in[2.31,2.61]\)
    (inference), below Liouville's \(10/3\).  As McKinnon--Roth note (ibid.\
    p.~24), at a \(k\)-rational point Liouville \(\alpha_x\ge\varepsilon_x\) is
    ``a stronger inequality than \(\alpha_x\ge\frac12\varepsilon_x\)''
    (\rec{papers_check/README.md} \S2.3).
  \item So this is the McKinnon--Roth regime.  According to the killer, the
    only predictions beyond Liouville/Seshadri are conjectural: Vojta, and
    McKinnon's conjecture that best approximations lie on rational curves.
    The killer also knew of no unconditional theorem that beats Liouville at a
    rational point of such a variety on a Zariski-dense set.  Round~22 checked
    this at snippet level (two independent searches plus an offline arXiv
    mirror): no unconditional theorem with \(\alpha^{\rm ess}>\tau_x\) at a
    rational point was found for any variety (\rec{round22/README.md} \S2,
    J3).  Seven full texts were then read (McKinnon--Roth 2015,
    2016~\cite{McKinnonRoth2015,McKinnonRoth2016}; the McKinnon--Satriano
    erratum~\cite{McKinnonSatriano2021}; Lehmann--McKinnon--Satriano
    2024~\cite{LMS2024}; Huang ANT 2021~\cite{Huang2021}, Acta Arith.\ 2019,
    Bull.\ SMF 2020).  Every unconditional lower bound on \(\alpha^{\rm ess}\)
    in them is Liouville with a section, at most \(\tau_x\).  Every exact value
    of \(\alpha^{\rm ess}\) they compute equals \(\tau\), with \(\tau\ge\dim\)
    (\rec{papers_check/README.md}).
\end{itemize}

\subsection{The approximation-constant form at the identity}

The following is from round~22 (\rec{round22/README.md} \S\S0, 4;
\rec{round22/toric.md}, \rec{round22/lit.md} with their killer reports) and
round~23.  The \(4\)-fold is the twisted Losev--Manin space \(X=\Ltw5\).
Propositions~\ref{prop:weakFano} and~\ref{prop:seshadri} are labelled
\emph{proved} in the record's summaries.  However, \rec{round22/README.md}
states that nothing written in round~22 was refereed by anyone other than the
killers, so these two results have been checked by the killers only.

\begin{proposition}[proved, killers only; round 22, \rec{round22/README.md} \S\S0, 4, \rec{round22/toric.md}]\label{prop:weakFano}
\(X\) is smooth and \emph{weak Fano}, not Fano: \(-K\) is nef and big but not
ample, with polytope the \(A_4\) root polytope, and \(-K=\phi^*(-K_{X'})\) for
the small crepant contraction \(\phi\) onto the twisted Gorenstein toric Fano
\(X'\), with \(X(\Q)=X'(\Q)=T_5(\Q)\) and \((-K)^4=70\).
\end{proposition}

With \(H=H_{-K}=R''^2\), an A-point at \(\delta\) has McKinnon--Roth exponent
\(\gamma=4/(1+2\delta)\), and the following hold.

\begin{proposition}[proved, killers only; round 22, \rec{round22/README.md} \S\S0, 4]\label{prop:seshadri}
\(\varepsilon_e(-K)=\alpha_e(-K)=2\), attained only on diagonal conics, which
carry repeated points.
\end{proposition}

\begin{theorem}[proved, refereed; \rec{round23/width.md} Thm 2.5]\label{thm:tau}
\(\tau_e(-K)=10/3\).  For every \(m\), the maximal order at \(e\) of a section
of \(m(-K)\) is \(\floor{10m/3}\), for the split and the twisted form, and it is
attained by Vandermonde products.
\end{theorem}

The proof slices the lattice points of \(a\Delta-b\Delta\) coordinatewise,
which gives the same family again, and peels the slices.  Consequently every
finite family of auxiliary Laurent polynomials over \(\Q(i)\), used through
Liouville with toric forced divisibility, gives for arbitrary profiles exactly
the Cilleruelo--C\'ordoba exponent \(\delta=1/10\) and never more; for
balanced profiles it gives \(\delta=1/30\).  McKinnon--Roth's base-locus
Liouville at \(e\) is therefore exactly Cilleruelo--C\'ordoba.
(McKinnon--Roth~\cite{McKinnonRoth2016}, Thm~3.3.  It needs \(L\) ample, so it
is applied to \(-K_{X'}\) on \(X'\).  The values are unchanged.)  The open window
\(\alpha^{\rm ess}\in(15/4,4)\) is untouched.

\(A_5(\delta<1/30)\) is exactly \(\alpha^{\rm ess}_e(-K)>15/4\) off a proper
closed set, and \((L_5^\circ)\) is the same over balanced sequences; Vojta
predicts \(\ge4=\dim X\).

\begin{named}{Question B, literature form}
Is the identity \(e\) \emph{essentially bounded} on the smooth weak Fano
\(4\)-fold \(X=\Ltw5\) (so that \(X(\Q)=T_5(\Q)\) for the anisotropic norm-one
torus \(T_5\)), in the sense of Lehmann--McKinnon--Satriano~\cite{LMS2024}
(Def~1.3: \(\alpha_{\rm ess}(e,-K_X+\varepsilon A)\ge\dim X\) for every ample
\(A\) and rational \(\varepsilon>0\))?  Since \(-K_X\) is big, this is
equivalent to \(\alpha^{\rm ess}_e(-K)\ge4\) (inference); the same constant is
obtained on \(X'\), since \(\phi\) is an isomorphism near \(e\).  Less
ambitiously: is
\(\alpha^{\rm ess}_e(-K)>15/4\), at least along sequences whose \(15\) toric
heights are asymptotically equal?
\end{named}

A yes gives the uniform Cilleruelo--Granville bound (\rec{round20/gauss.md}
Thm~3.4\('\)).  Vojta's Main Conjecture implies it.  McKinnon--Satriano's
erratum~\cite{McKinnonSatriano2021} (correcting their Prop~1.3) has Prop~0.1:
``Let \(X\) be a smooth projective variety over a number field \(k\), and let
\(P\in X(k)\) \dots\ If \(-K_X\) is big, then Vojta's Main Conjecture implies
that \(X\) is canonically bounded at \(P\).''  Lehmann--McKinnon--Satriano
Prop~1.4 gives the same for essential boundedness on any smooth projective
\(X\).  Neither assumes splitness.  Both use Vojta on \(X\) with a
normal-crossings divisor of multiplicity \(\dim X\) through \(P\), a different
instance from \(V_E(5)\) with the same conclusion.  In
Lehmann--McKinnon--Satriano Thm~1.5 essential boundedness is a hypothesis.  No
unconditional proof was found for any variety where \(\tau_P(-K)<\dim X\).
Full texts were checked for seven papers; McKinnon--Satriano (TAMS) and
Monahan--Satriano remain unread (\rec{papers_check/README.md}).

\subsection{The split calibration and \texorpdfstring{Erd\H{o}s \#887}{Erdős \#887}}

The following is from round~23 (\rec{round23/split.md}, refereed).  The same
window occurs on the \emph{split} Losev--Manin \(4\)-fold at \(e\).  There a
primitive positive \(5\)-tuple \(u\) has
\(H_{-K}=(\max u/\min u)\,\mathrm{lcm}(u)\) and
\(\dist\asymp(\max u-\min u)/\min u\).

\begin{named}{Split Question B (not known)}
Do five distinct integers with gcd \(1\) in
\([U,U(1+\mathrm{lcm}^{-(1+2\delta)/4})]\), \(\delta<1/30\), have their ratio
class in a fixed proper closed set once the lcm is large?  Equivalently, this
concerns five factorisations of \(N\) within hyperbolic angle
\(N^{-(1+2\delta)/4}\) on \(xy=N\).
\end{named}

\begin{theorem}[proved, refereed; \rec{round23/split.md}]\label{thm:split}
Split Question~B implies Erd\H{o}s Problem
\#887 in its absolute form (\(\exists K\,\forall C\): eventually at most \(K\)
divisors of \(n\) in \([X,X(1+Cn^{-1/4})]\), for every \(X\)), and Ruzsa's
\#886 for \(\varepsilon>(15/64)(1+2\delta)\).
\end{theorem}

Split Question~B is \emph{not known}, and both Erd\H{o}s problems are open
(formal-conjectures, October 2026).  The proof of the implication uses a
random \(5\)-subset, averaging over threshold layers, and Schwartz--Zippel.
The 2026 five-divisor construction for \#887
(\(n=(x^2-1)(x^2-4)(x^2-16)(x^2-49)\), \(z^2=56x^2+65\)) is a rational curve
hitting each of the \(15\) cuts once, with \(\alpha=15/4\) exactly.  It lies in
a \(2\)-parameter family sweeping a hypersurface \(W\), which every admissible
split exceptional set must contain.  Such curves do not exist on the twisted
form: for a curve over \(\Q\) there, \(E_A\cdot C=E_{A^c}\cdot C\), so a
balanced twisted curve at \(15/4\) has \(B\cdot C\ge30\) and
\(\mult_e\ge8\).  A proof of Question~B valid for every \(\Q\)-form of the
\(4\)-fold would therefore solve \#887; a proof for \(T_5\) alone may need the
anisotropy.  Split Question~B is undecided.  It would fail if infinitely many
such hypersurfaces existed, or if absolute \#887 failed.

\subsection{Anisotropy}

The following results are from round~24 (\rec{round24/theory.md},
\rec{round24/search.md}, refereed).

\begin{theorem}[refereed; \rec{round24/theory.md}; proved except where labelled, \rec{round24/README.md} \S1]\label{thm:aniso-local}
Pointwise, the anisotropy does nothing.  Let \(\LL\) be the arguments
``archimedean Taylor upper bound \(\ge\) product-formula lower bound'' whose
lower bound may use integrality, toric forced divisibility at every place of
\(\Q(i)\), point parity \(v_{\conj\pi}=-v_\pi\) and reality of values, with
auxiliary objects that are isotypic sections for
\(\sigma_X=\iota\circ\overline{\phantom{x}}\), polynomials in group words of
A-points (including \(\iota y\)), or polynomials in the Hilbert-90 vectors
(\(x_j=P_j/\conj P_j\)).  Every such inequality is implied by toric Liouville
for one section (or its \(n\)-point form at the generating points).  So on
balanced simple profiles with disjoint primes nothing in \(\LL\) beats
\(15/4\) (\(\delta=1/30\)), and nothing beats \(\delta=1/10\) for arbitrary
profiles, for every family with
\(\log H(F)+\sum d_i\log d_i=o(\sum d_iw_i)\).
\end{theorem}

The reason is that \(\sigma_X\) fixes \(e\) and orders at \(e\), so Galois
averaging is lossless; the isotypic parts of \(H^0(m(-K))\) have maximal orders
\(\floor{10m/3}\) and \(\floor{10m/3}-1\) in both forms.  Twist-sensitive
congruence divisibility at inert places contributes \(O_F(1)\).

\begin{theorem}[refereed, component claim numerical; \rec{round24/theory.md}, \rec{round24/search.md}, \rec{round24/README.md} \S1.2]\label{thm:aniso-global}
Globally, the anisotropy is visible only through curves.  By the parity lemma
a balanced twisted curve at \(15/4\) has \(B\cdot C\ge30\).  Such curves exist
over \(\R\): an explicit one-parameter real family \(M\) (two conjugate
branches of contact \(4\), each ray hit once).  Only finitely many curves of
its complex component are defined over \(\Q\) (geometric monodromy \(S_4\) of
the place configuration, Faltings; ineffective).  The same finiteness holds
for the split form, so this is not an anisotropy effect.
\end{theorem}

That \(M\) is a component rests on a numerical computation: the Jacobian has
rank \(28\) along \(M\) (\rec{round24/README.md} \S1.2, where the component
claim is labelled \emph{numerical}).  The finiteness is ineffective; a finite
search found no such curve over \(\Q\) of height \(\le300\).

Whether the class has other components is open.  Blind numerical solvers all
missed \(M\), so failed searches in this class are not evidence of emptiness.
Consequently a proof of Question~B for \(T_5\) that is not also a proof for
the split form (hence of Erd\H{o}s \#887) must use a global twist-sensitive
input; the only one known is the parity lemma.

\subsection{The absolute form of the A-branch}

\begin{theorem}[proved; round 23 judge J4, independent referee \rec{round23/referee\_J4.md}]\label{thm:absolute}
If \(A_5(\delta)\) holds for one \(\delta\in(0,1/30)\), there is
\(K=K(\delta,Z)\) such that for every \(C>0\) and every
\(R\ge R_1(C)=\max(C^{2/\kappa},(CR_0)^2)\), every arc of length
\(\le C\sqrt R\) holds at most \(K\) lattice points; so
\(\sup_C M_\infty(C)<\infty\).  Here \(\kappa=1-(15/16+\eta)(1+2\delta)>0\)
and \(\eta=(1/(1+2\delta)-15/16)/2\).  The same argument gives this from
\(A_k(\delta)\), \(\delta<1/(2^k-2)\), for every \(k\ge3\): convexity averaging
of the tuple gcd alone reaches the sharp range, with no balanced selection and
no unit-class restriction.
\end{theorem}

Scope: this is proved from the \(A_5\) branch only, not from \((L_5^\circ)\)
(whose route subdivides the arc and gives \(K\) linear in \(C\)).  So a
construction with \(M_\infty(C)\to\infty\) would refute \(A_5(\delta)\) for all
\(\delta<1/30\), and \(V_E(5)\), but not \((L_5^\circ)\).  The absolute form is
only formally stronger than \(M(C)<\infty\), and it also follows from
Cilleruelo--Granville's own \(R^{1-\varepsilon}\) conjecture
(\(\varepsilon=1/4\)).

\begin{question}\label{q:B}
Prove \((L_5^\circ)\), or \(A_5(\delta)\) for some \(\delta<1/30\).  Or refute
\((L_5^\circ)\) by a Zariski-dense supply of balanced bounded-residue
\(5\)-tuples at exponent \(\ge7/15\) (for a rational family, one dominating
the \(4\)-dimensional \(T_5\)); that would refute \(V_E(5)\).  Unlike \(k=4\),
a \(2\)-parameter family plus its curve parameter has image of dimension
\(\le3\) and refutes nothing.  By the killer's dimension count, one-variable
polynomial-block families cannot be dense at \(k=5\) (\emph{heuristic}), and
bounded-degree covering families are excluded (Theorem~\ref{thm:covering},
that is, \rec{round20/gauss.md} Thm~4.3(ii)) (\rec{round20/kill\_gauss.md}
\S\S4, 6.3).  The two-branch twisted Pell ansatz (\(B\cdot C=30\), contact
\(4\), round~24) yields real curves at exactly \(15/4\) but only finitely many
over \(\Q\) on the known component; it can only enlarge \(Z\).
\end{question}

\paragraph{Compared with \((L_8)\)} (\rec{round20/README.md} \S\S5, 7).
\((L_5^\circ)\) is smaller: a \(4\)-fold and \(5\) points, against a
\(5\)-fold and \(8\) points.  It needs the archimedean divisor \(D=E\) rather
than \(D=0\), and its heuristic margin is much thinner:
\(\kappa=2\delta\approx1/15\) and \(e^{-w}\) per scale, against
\(K\cdot\beta=10\) and \(e^{-10w}\).  Neither statement is known to imply the
other.  Fakes violate \((L_5^\circ)\), whereas \((L_8)\) is fake-blind.

\section{Question C: positive-gain non-unit pencils on \texorpdfstring{\(\Mbar\)}{M0,8}}\label{sec:C}

This question is smaller than Questions~\ref{q:A} and~\ref{q:B}.

\subsection{Gain and the abc criterion}

The following is from \rec{round18/theory.md} \S2.4 and
\rec{round19/theory.md} \S1.  Let \(X=\Mbar\), let \(\partial=\sum_TD_T\)
(\(119\) components, snc), and let \(\beta\) be the balanced class.  For a
pencil (a non-constant rational function) \(f\), write \(F_p=f^*(p)\sim L\) and
\(R_p=(F_p)_{\rm red}\cdot\beta\).

\begin{definition}[gain]
\[
  \gain(f)=L\cdot\beta-(R_0+R_1+R_\infty)
  =\sum_{Y\subset F_0+F_1+F_\infty}\tfrac{m_Y-3}{3}\,Y\cdot\beta .
\]
\end{definition}

When the base locus lies in codimension-\(2\) boundary strata, this equals
\(10-K^{\log}_{\tilde X/\PP^1}\cdot\tilde\beta\).

A \emph{unit pencil} is \(f=\lambda A/C\) with \(\lambda\in\Q^*\), where
\(A,C\) are coprime products of differences \(x_i-x_j\) in Kapranov's chart.
Put \(B=C-\lambda A\).

\begin{proposition}[conditional on abc, implication proved; \rec{round18/theory.md} Prop 2.6]\label{prop:abc}
Assume the abc conjecture.  If \(f\) is a unit pencil with \(\gain(f)>0\), then
\((L_8^{pc})\) holds with \(Z=\{B=0\}\).
\end{proposition}

The mechanism is as follows.  On \(\Bal^{pc}\) we have
\(\gcd(A(x),C(x))=\prod d_T^{\mu_T}\), so
\(h(f(\Phi))=(L\cdot\beta)w+O(\eta w)\).  Divide \((\lambda A(x),B(x),C(x))\) by
its gcd to get a coprime abc triple.  The log of the radical of that triple is
\(\le(F_{\rm red}\cdot\beta)w+O(\eta w)\).  The refereed abc barrier covers
only relations among products of differences (\rec{round8/weakest.md}
Prop~4.2), so this was a genuine gap in the barrier map.

\subsection{Round 19 closed the unit part}

The following is from \rec{round19/theory.md}, refereed in
\rec{round19/referee\_theory.md}; a summary is in \rec{round19/README.md}
\S1.

\begin{theorem}[proved; \(\beta\)-weighted log Riemann--Hurwitz; \rec{round19/theory.md} Thm 2.4]\label{thm:RH}
For every non-constant rational function \(f\) on \(\Mbar\),
\[
  \sum_{p\in\PP^1}\big(L\cdot\beta-R_p^{\rm int}\big)
  \ \le\ 2L\cdot\beta+\Lambda_{k(f)}\cdot\beta
  \ \le\ 2L\cdot\beta+31 .
\]
Here:
\begin{itemize}
  \item \(R_p^{\rm int}\) is the \(\beta\)-degree of the reduced non-boundary
    part of \(F_p\);
  \item \(k(f)\) is the least \(|S|\) such that \(f\) factors through a
    forgetful map \(\pi_S\);
  \item \(\Lambda_k=\pi_k^*\psi_k\), with \(\Lambda_k\cdot\beta=16,24,28,30,31\)
    for \(k=4,\dots,8\).
\end{itemize}
Equality holds for the cross-ratio and for \(4z(1-z)\).
\end{theorem}

\begin{corollary}[proved; \rec{round19/theory.md} Cor 2.5]\label{cor:gain31}
\(\gain(f)\le31-n'(f)\), where \(n'(f)=\#\{T: f|_{D_T}\text{ constant}\}\).
\end{corollary}

\paragraph{Idea of the proof.}
The form \(df\) spans a line subsheaf of \(\Omega^1_X(\log\partial)\), and its
divisor is read off the fibres.  The forgetful filtration
\(F_k=\pi_k^*\Omega^1_{\overline M_{0,k}}(\log\partial_k)\) has line-bundle
quotients \(\Lambda_k\).  Every line subsheaf of \(F_k\) that is not in
\(F_{k-1}\) has \(\beta\)-degree \(\le\Lambda_k\cdot\beta\), because \(\beta\) is
positive on every effective divisor (Keel--McKernan, via
\rec{threshold/L8.md} Thm~A(a)).  This is not a curve argument.  Restricting to
covering curves of class \(d\beta\) only gives gain \(\le(2g-2)/d\), a bound
that is \(\ge10\) and so cannot exclude positive gain.

\begin{theorem}[proved; \rec{round19/theory.md} Thm 3.3, Cor 3.4]\label{thm:units}
For every non-constant unit \(u\) of \(M_{0,8}\) and every
\(\lambda\in\C^*\),
\[
  \gain(\lambda u)\le\Lambda_{k(u)}\cdot\beta-N(u)-n_1\le-2^{8-k(u)}\le-1
  \qquad(\le-16\ \text{if}\ k(u)=4),
\]
where \(n_p=\#\{T:D_T\subset\operatorname{supp}F_p\}\) counts the boundary
divisors in the fibre over \(p\).  The same bound \(\gain\le-1\) holds for
every composite \(h\circ u\).
\end{theorem}

The key input is \(N(u)=\#\{T:\ord_{D_T}u\ne0\}\ge32\), from multilinear
Schwartz--Zippel on \(\{0,1\}^7\).  So Proposition~\ref{prop:abc} never
applies.  Exact computation agrees: gain \(\le-32\), attained, for all unit
pencils of degree \(D\le6\) and all \(\lambda\in\Q^*\) (\rec{round19/search.md},
reproduced by the referee; \rec{round18/theory.md} \S2.4).  Morphisms have
gain \(\le-32\) (Bruno--Mella, \emph{read from the abstract only}).

\subsection{What is open}

The following is from \rec{round19/theory.md} \S4 and
\rec{round19/referee\_theory.md} m2.  The remaining case is non-unit pencils
in
\[
  P^\circ=\Big\{f\ \text{regular on}\ X\setminus\textstyle\bigcup_{T\ne T'}(D_T\cap D_{T'})\Big\}.
\]
Two things are needed, and both are open.
\begin{enumerate}[label=(\roman*),leftmargin=2em]
  \item \emph{A gcd extension of Proposition~\ref{prop:abc}} to \(P^\circ\).
    This means bounding the excess of \(\gcd(A(x),C(x))\) over
    \(\prod d_T^{\mu_T}\) at balanced pairwise-coprime points.  Round~18
    asserted, without proof, that for base locus in codimension-\(2\) boundary
    strata the excess gcd vanishes on \(\Bal^{pc}\) (\rec{round18/theory.md}
    \S2.4, after Prop~2.6).  Round~19 \S4.1 records this as plausible but
    \emph{unproved} and gives only an unrefereed sketch of the obstruction.
    (For base locus meeting the interior, outside \(P^\circ\), round~18 calls
    the gcd ``a Vojta-type gcd problem in its own right''.)
  \item \emph{A pencil with gain \(\ge1\).}  Any such pencil must satisfy these
    proved necessary conditions:
    \begin{itemize}
      \item \(n'(f)\le31-2^{8-k(f)}\) and \(k(f)\ge5\);
      \item \(f\) is not of the form \(h\circ u\) with \(u\) a unit;
      \item all three fibres are non-reduced, with \(\beta\)-average
        multiplicities \(\mu_p=L\cdot\beta/R_p\) satisfying
        \(\sum_{p\in\{0,1,\infty\}}1/\mu_p<1\);
      \item \(\sum_{p\notin\{0,1,\infty\}}(L\cdot\beta-R_p^{\rm int})\le30-n(f)\),
        where \(n(f)=n_0+n_1+n_\infty\) counts the boundary divisors in the
        three special fibres (\(n(f)\le n'(f)\)).
    \end{itemize}
    \emph{No example is known}, and Theorem~\ref{thm:RH} alone allows gain
    \(\le31\).  Campana-type statements see only gcd multiplicities, so they
    do not decide the question.
\end{enumerate}

Even with both (i) and (ii), the result would be only an abc-conditional
\((L_8^{pc})\).  Reaching \((L_8)\) also needs a nested-chain extension.
There, for the chain \(\{1,2\}\subset\{1,2,3\}\), the gcd exceeds the value in
Proposition~\ref{prop:abc} by a full block (\rec{round18/README.md} \S3.3, J2;
\emph{unrefereed}).  The judge adds that the radical also counts \(p\) once,
which might compensate, and that this is not analysed; so J2 shows a gcd
excess, not that the extension fails.

\begin{question}\label{q:C}\leavevmode
\begin{enumerate}[label=(\alph*)]
  \item Does some \(f\in P^\circ\) have
    \(K^{\log}_{\tilde X/\PP^1}\cdot\tilde\beta<10\), that is,
    \(\gain(f)>0\)?
  \item Does Proposition~\ref{prop:abc} (\rec{round18/theory.md} Prop~2.6)
    extend to \(P^\circ\)?  This would be a gcd inequality for two forms at
    balanced points.
  \item Is Theorem~\ref{thm:RH} (\rec{round19/theory.md} Thm~2.4) known, or
    useful elsewhere?  The record notes that it holds verbatim for any movable
    class \(\gamma\), with \(31\) replaced by \(\max_k\Lambda_k\cdot\gamma\),
    the maximum taken over orderings (unused there; \rec{round19/theory.md}
    \S7 item~4, referee m7).  Is it known, or sharp, in that generality, and
    on \(\overline M_{0,n}\)?
\end{enumerate}
\end{question}

\section{What any proof must avoid}\label{sec:avoid}

Each item below is a barrier result in the record
(\rec{overview/overview.md} \S\S5.4, 7.6).  Statuses are as given there.
\begin{enumerate}[leftmargin=2em]
  \item \emph{Local inputs.}  These are pairs, gcds, cuts, monomial
    (Liouville) gaps, rational characters, and point-level residues at moduli
    \(\le M^A\) with \(A\) fixed in advance.  Fakes at \(W=O(M\log M)\) satisfy
    all of them (\emph{proved}, with the scope limits of
    \rec{overview/overview.md} \S5.2 item~2).  Caveat: caps uniform in \(A\)
    are not proved, and P3-all\(^*\) is uncapped.
  \item \emph{Averaged or robust analytic inputs.}  This covers large sieves
    with non-sharp constants, prime sums or box counts with power-saving
    error, and ``almost all'' statements.  Grafted fakes satisfy them
    (\rec{round10/analytic.md} Thm~B, \emph{proved}; composite-supported
    version \rec{round20/composite.md}, killer-checked).  Exceptions inside the
    low-height window are listed in \rec{round21/README.md}.
  \item \emph{Auxiliary polynomials and single-place Diophantine
    approximation.}
    \begin{itemize}
      \item On \(X_M\): \(\tau^*(M)=2-1/\ceil{M/2}<2\), so Liouville at
        \(\infty\), and Ru--Vojta with \(S=\{\infty\}\), never reach exponent
        \(\frac12\) (\rec{threshold/upper.md}, \emph{proved}).
      \item On \(\overline M_{0,k}\) for any \(k\): \(\beta\) lies in the
        interior of the movable cone (\rec{threshold/L8.md} Thm~A,
        \emph{proved} as a formal barrier).
      \item On the Gaussian model: the forms are exactly critical
        (Lemma~\ref{lem:K1} with \rec{threshold/construct.md} Cor~C3 and
        Thm~B).
      \item Roth- and Schmidt-type bounds (McKinnon--Roth) give
        \(\alpha^{\rm ess}\ge\beta<\tau\) at a rational target: weaker than
        Liouville.
    \end{itemize}
  \item \emph{Single linear forms in logarithms} (Baker--Matveev, single-form
    Lang--Waldschmidt).  Fakes satisfy them (\rec{round8/lw.md} Thms~A, B;
    \rec{round8/stability.md} Prop~6.1; \emph{proved}).
  \item \emph{\(S\)-unit and Subspace arguments with thresholds depending on
    \(S\) only through \(|S|\).}  These give at best \(M\le\omega(N')\) (lit18
    Props~A, B; \emph{unrefereed}).  Fixed-\(S\) theorems such as
    \rec{round18/theory.md} Thm~2.5 have \(S\)-dependent constants, and \(S\) is
    the unbounded set of primes of \(N\).
  \item \emph{abc on single relations.}  abc through relations among products
    of chords (pair, triangle, Ptolemy or \(n\)-term) is vacuous on squarefree
    profiles (\rec{round8/weakest.md} Prop~4.2), which contain every recorded
    fake (Cor~4.3), so it cannot prove growth (\emph{proved}).  abc on unit
    pencils is closed as well (Theorem~\ref{thm:units}).
  \item \emph{gcd normalisation.}  It is invisible to forced-divisibility,
    auxiliary-form, Liouville, pair, cut, character and single-form
    inequalities (Lemma~\ref{lem:K1}, proved by the killer).  Vojta's \(D=E\)
    term and the heuristic count do see it.
  \item \emph{Genus-\(0\) Arakelov theory and unlikely intersections in tori.}
    The first reduces to items~1 and~3 (\rec{overview/overview.md} \S5.4 rows~2
    and~8) and to the limitation ``Lim'' of known tools (Mordell--Weil packing
    gives \(e^{O(\mathrm{rank})}\); bounded rank is equivalent to the uniform
    theorem).  The second is vacuous, or capped at
    \(M\le\omega(N')+O_C(1)\); the R\'emond-cone case of that ceiling is only
    \emph{heuristic} (\rec{round20/arakelov.md}, \rec{round20/wild.md};
    killer-checked; \rec{overview/overview.md} \S5.4 rows~18--19).
  \item \emph{For refutations.}
    \begin{itemize}
      \item No feasible computation exists (\rec{round18/census.md} Thm~E).
      \item No density mechanism works: proved at \(k=8\), heuristic at
        \(k=5\) for circles.
      \item Bounded-degree covering curve families whose members each carry
        infinitely many balanced small-\(H\) points are excluded at \(k=8\)
        (genus \(\le1\) by \rec{round18/census.md} Prop~7.5; genus \(\ge2\) by
        Faltings), and, for circles, those whose members each carry infinitely
        many A-points are excluded at \(k\ge5\), \(\delta>0\)
        (Theorem~\ref{thm:covering}).  Families with infinitely many rational
        points per member are not excluded as such (the lines of Kapranov's
        \(\PP^5\) are one; correction~M1).
      \item What remains are dense supplies with finitely many points per
        curve, curves of unbounded degree, and sporadic points.
    \end{itemize}
\end{enumerate}

\section{Where the details are}\label{sec:pointers}

All files are in the repository \rec{github.com/craigxchen/Jarnik}, branch
\rec{claude/upbeat-einstein-p0peju}, directory \rec{docs/claude\_review/}.
Each write-up \rec{X.md} has its referee report \rec{referee\_X.md} next to it,
or, in round~20, a kill report \rec{kill\_X.md}.  Its exact checks are in
\rec{X\_checks/}.  The round READMEs give the referee-corrected statements.

\begin{center}
\small
\begin{tabular}{@{}p{0.40\textwidth}p{0.55\textwidth}@{}}
\hline
\textbf{Item} & \textbf{File} \\
\hline
consolidated account (status labels, corrections) &
  \rec{overview/overview.md} (\S\S4, 5.4--5.5, 7.1, 7.1A, 7.6, 8.2) \\
Theorem~\ref{thm:twothirds}, the \((\frac23+\varepsilon)\) bound &
  \rec{two\_thirds/proof.md} (Thm~7.5, Cor~8.1, \S10) \\
Vojta on \(\mathrm{Bl}_YX_7\) &
  \rec{growth/conditional.md} (Thms~5.1--5.2),
  \rec{growth/referee\_conditional.md} \\
\((L_8)\); Vojta on \(\Mbar\) implies uniform &
  \rec{growth/ptolemy.md} (\S6 Thm~D, \S7), \rec{growth/referee\_ptolemy.md} \\
\((L_5)\), \((L_6)\) false; Liouville barrier on \(\overline M_{0,k}\) &
  \rec{threshold/L8.md} (Thms~A, B, Prop~C), \rec{threshold/referee\_L8.md},
  \rec{growth/referee\_ptolemy.md} \S3.3 \\
round 18: census, sandwich, lattice form, lifts, Prop~2.6 &
  \rec{round18/README.md}, \rec{round18/theory.md}, \rec{round18/census.md},
  \rec{round18/lift.md} (with \rec{referee\_*.md}) \\
round 19: Thm~2.4, Thm~3.3, non-unit residue &
  \rec{round19/README.md}, \rec{round19/theory.md}, \rec{round19/search.md}
  (with \rec{referee\_*.md}) \\
round 20: \((L_5^\circ)\), \(V_E(5)\), Thm~4.1, Lemma~K1 &
  \rec{round20/README.md}, \rec{round20/gauss.md}, \rec{round20/kill\_gauss.md} \\
round 21: low-height window (open) & \rec{round21/README.md} \\
round 22: approximation constants at \(e\) &
  \rec{round22/README.md}, \rec{round22/toric.md}, \rec{round22/lit.md} \\
round 23: \(\tau_e(-K)=10/3\), split calibration, absolute form &
  \rec{round23/width.md}, \rec{round23/split.md}, \rec{round23/referee\_J4.md} \\
round 24: anisotropy & \rec{round24/theory.md}, \rec{round24/search.md} \\
fake circles and capped classes &
  \rec{round4/outside.md}, \rec{round5/construct.md}, \rec{round8/lw.md},
  \rec{round10/analytic.md} \\
mean-width Thm~B, Cor~C3 (auxiliary forms capped) &
  \rec{threshold/construct.md} \\
\hline
\end{tabular}
\end{center}

\begin{thebibliography}{9}

\bibitem{CillerueloCordoba1992}
J. Cilleruelo and A. Cordoba,
\newblock Trigonometric polynomials and lattice points,
\newblock \emph{Proceedings of the American Mathematical Society}
\textbf{115} (1992), no. 4, 899--905.

\bibitem{CillerueloGranville2007}
J. Cilleruelo and A. Granville,
\newblock Lattice points on circles, squares in arithmetic progressions and
sumsets of squares,
\newblock in \emph{Additive Combinatorics}, CRM Proceedings \& Lecture Notes
\textbf{43}, American Mathematical Society, Providence, RI, 2007, 241--262.

\bibitem{CillerueloGranville2009}
J. Cilleruelo and A. Granville,
\newblock Close lattice points on circles,
\newblock \emph{Canadian Journal of Mathematics} \textbf{61} (2009), no. 6,
1214--1238.

\bibitem{Huang2021}
Z. Huang,
\newblock Rational approximations on toric varieties,
\newblock \emph{Algebra \& Number Theory} \textbf{15} (2021), no. 2, 461--512.

\bibitem{LMS2024}
B. Lehmann, D. McKinnon and M. Satriano,
\newblock Approximating rational points on surfaces,
\newblock arXiv:2403.02480 (2024).

\bibitem{McKinnonRoth2015}
D. McKinnon and M. Roth,
\newblock Seshadri constants, diophantine approximation, and Roth's theorem
for arbitrary varieties,
\newblock \emph{Inventiones Mathematicae} \textbf{200} (2015), 513--583.

\bibitem{McKinnonRoth2016}
D. McKinnon and M. Roth,
\newblock An analogue of Liouville's theorem and an application to cubic
surfaces,
\newblock \emph{European Journal of Mathematics} \textbf{2} (2016), no. 4,
929--959.

\bibitem{McKinnonSatriano2021}
D. McKinnon and M. Satriano,
\newblock Approximating rational points on toric varieties,
\newblock \emph{Transactions of the American Mathematical Society}
\textbf{374} (2021), no. 5, 3557--3577; erratum (2023).

\end{thebibliography}

\end{document}
